\documentclass[11pt,reqno]{amsart}
\usepackage[margin=1in]{geometry}
\usepackage{amsmath,amssymb,amsthm,mathtools}
\usepackage{booktabs}
\usepackage{microtype}
\usepackage[colorlinks=true,linkcolor=blue,citecolor=blue,urlcolor=blue]{hyperref}

\theoremstyle{plain}
\newtheorem{theorem}{Theorem}[section]
\newtheorem{proposition}[theorem]{Proposition}
\newtheorem{corollary}[theorem]{Corollary}
\newtheorem{lemma}[theorem]{Lemma}
\newtheorem{conjecture}[theorem]{Conjecture}
\theoremstyle{definition}
\newtheorem{definition}[theorem]{Definition}
\theoremstyle{remark}
\newtheorem{remark}[theorem]{Remark}

\newcommand{\R}{\mathbb{R}}
\newcommand{\C}{\mathbb{C}}
\newcommand{\F}{\mathbb{F}}
\newcommand{\N}{\mathbb{N}}
\newcommand{\PP}{\mathbb{P}}
\newcommand{\EE}{\mathbb{E}}
\newcommand{\one}{\mathbf{1}}
\newcommand{\smin}{\sigma_{\min}}
\newcommand{\dF}{d_{\F}}
\DeclareMathOperator{\Beta}{Beta}
\DeclareMathOperator{\Var}{Var}

\title[The sharp exponent for coordinate-optimized least singular values]{The sharp exponent for the smallest singular value of a coordinate-optimized Gaussian submatrix}
\author{Research draft}
\date{July 6, 2026}

\begin{document}
\begin{abstract}
Let $A$ be an $N\times N$ matrix with independent standard real or complex Gaussian entries, write
$\dF\in\{1,2\}$ for the real dimension of the field, and for $\alpha\in(0,1)$ set
$n=\lfloor\alpha N\rfloor$ and
\[
  M_N(\alpha)=\min_{|S|=|T|=n}\smin\bigl(A[T,S]\bigr),
\]
the least value of the smallest singular value over all pairs of coordinate sets $S,T\subseteq[N]$
of size $n$, where $A[T,S]$ is the $n\times n$ submatrix on rows $T$ and columns $S$. We prove
\[
  -\frac1N\log M_N(\alpha)\ \xrightarrow{\ \PP\ }\ \frac{2H(\alpha)}{\dF},\qquad
  H(\alpha)=-\alpha\log\alpha-(1-\alpha)\log(1-\alpha).
\]
This is the sharp large-deviation exponent for the worst conditioning achievable by restricting a
Gaussian linear map to an $\alpha$-fraction of its coordinates in domain and range. The value settles,
for square submatrices, the two-sided problem for which only the surrounding bounds $H(\alpha)/\dF$
and $2H(\alpha)/\dF$ were available. More generally, for submatrices with $r$ surplus rows the same
argument gives the exponent $2H(\alpha)/(\dF(r+1))$ for every fixed $r$. The upper bounds are
first-moment estimates. The lower bounds are second-moment arguments on a marked-row construction;
the new ingredient is a small-ball estimate for two blocks of linear functionals of Gaussian vectors
taken along two directions of arbitrary overlap, sharp up to a logarithmic factor uniformly in the
overlap, which removes the loss a crude bound incurs on heavily overlapping submatrices. Two
consequences are recorded: the achievable rate for support-optimized products of Gaussian matrices
improves to $2H(\alpha)/\dF$, and for Rademacher entries the minimum vanishes for all
$\alpha\le\tfrac12$, so the Gaussian constant is not distribution-free.
\end{abstract}
\maketitle

\section{Introduction}\label{sec:intro}
Throughout $\F\in\{\R,\C\}$, $\dF=\dim_\R\F\in\{1,2\}$, and $A=(a_{ij})_{i,j=1}^N$ has independent
entries, each a standard $\F$-Gaussian: $a_{ij}\sim N(0,1)$ when $\F=\R$, and $a_{ij}$ a standard
complex Gaussian with $\EE|a_{ij}|^2=1$ when $\F=\C$. For $S,T\subseteq[N]$ we write $A[T,S]$ for the
submatrix with rows indexed by $T$ and columns by $S$. Fix $\alpha\in(0,1)$, put $n=\lfloor\alpha
N\rfloor$, and study
\begin{equation}\label{eq:MN}
  M_N(\alpha)=\min_{\substack{S,T\subseteq[N]\\ |S|=|T|=n}}\smin\bigl(A[T,S]\bigr).
\end{equation}
Because $A$ has independent entries, each fixed $A[T,S]$ is itself an $n\times n$ matrix with
independent standard $\F$-Gaussian entries, so its smallest singular value is polynomially small in
$n$ (the hard edge). The minimum in \eqref{eq:MN} is taken over $\binom Nn^2=e^{2H(\alpha)N+o(N)}$
strongly dependent submatrices, and it is exponentially small. Our result identifies the exponent.

\begin{theorem}\label{thm:main}
For every $\alpha\in(0,1)$,
\[
  -\frac1N\log M_N(\alpha)\ \xrightarrow{\ \PP\ }\ \frac{2H(\alpha)}{\dF}\qquad(N\to\infty).
\]
\end{theorem}

The same holds with surplus rows. For an integer $r\ge0$ put
\begin{equation}\label{eq:MNr}
  M_N^{(r)}(\alpha)=\min_{\substack{T,S\subseteq[N]\\ |T|=n+r,\ |S|=n}}\smin\bigl(A[T,S]\bigr),
\end{equation}
the minimum over all $(n+r)\times n$ coordinate submatrices; $M_N^{(0)}=M_N$.

\begin{theorem}\label{thm:surplus}
For every $\alpha\in(0,1)$ and every fixed integer $r\ge0$,
\[
  -\frac1N\log M_N^{(r)}(\alpha)\ \xrightarrow{\ \PP\ }\ \frac{2H(\alpha)}{\dF(r+1)}\qquad(N\to\infty).
\]
\end{theorem}

Section~\ref{sec:surplus} proves Theorem~\ref{thm:surplus}; the case $r=0$ recovers
Theorem~\ref{thm:main}, and Sections~\ref{sec:construct}--\ref{sec:second} present that case
separately because all the ideas are visible there with lighter bookkeeping. Section~\ref{sec:conseq}
records some consequences and complements: for support-optimized products, for Rademacher entries,
and for tall submatrices, where support optimization is powerless and the minimum stays of order
$\sqrt N$.

\paragraph{Relation to previous work.}
The one-sided statistic $\min_{|T|=m+r}\smin(G_T)$, where $G\in\F^{N\times m}$ is Gaussian and the
minimum is over row sets $T$ of size $m+r$ with columns fixed, was determined by Shmalo
\cite{shmalo}: for fixed surplus $r$ it satisfies $\tfrac{r+1}{m}\log(\cdot)\to-h(\gamma)/\dF$, where
$\gamma=\lim N/m$ and $h(\gamma)=\gamma\log\gamma-(\gamma-1)\log(\gamma-1)$. Specializing to a square
window ($r=0$, columns all of $[m]$) and then also minimizing over columns is exactly \eqref{eq:MN}.
For that two-sided minimum \cite{shmalo} records the bracket
\begin{equation}\label{eq:bracket}
  \frac{H(\alpha)}{\dF}\ \le\ \liminf_N-\tfrac1N\log M_N(\alpha)\ \le\ \limsup_N-\tfrac1N\log M_N(\alpha)\ \le\ \frac{2H(\alpha)}{\dF},
\end{equation}
with the sharp constant left open; the lower end used only a single fixed column set, while the
upper end is a union bound over all pairs. It was not known whether the exponent exceeds
$H(\alpha)/\dF$ at all, i.e.\ whether optimizing over columns helps at exponential order. Theorem
\ref{thm:main} gives the sharp value $2H(\alpha)/\dF$, the top of the bracket. We use the identity
$H(\alpha)=\alpha\,h(1/\alpha)$, which reconciles the two normalizations.

\paragraph{The point of the proof.}
The upper bound in \eqref{eq:bracket} is a first moment (Section~\ref{sec:upper}); the content is the
matching lower bound. The natural route is a second moment for the number of near-singular
submatrices, but the pairs of submatrices that overlap in almost all of their rows and columns are
too correlated for the crude estimate that suffices at the typical scale, and they are numerous
(this is the obstruction identified in \cite{shmalo}). Two devices resolve it. First, a marked-row
construction (Section~\ref{sec:construct}) reduces near-singularity of $A[T,S]$ to a single linear
constraint through one distinguished row; drawing that row from a small fixed set $M$ while barring
the remaining rows from $M$ forces two constructed events with distinct marked rows to be exactly
independent, so all correlation is carried by pairs sharing a marked row. Second, for such pairs the
correlation is governed by one overlap parameter, the inner product of the two constructed normals,
whose small-ball behaviour we control by a sharp two-direction estimate (Section~\ref{sec:twoslab})
together with the fact that the constructed normal is Haar-distributed. The crude bound on heavily
overlapping pairs is replaced by these, and the second moment concentrates
(Section~\ref{sec:second}).

\paragraph{An interpretation.}
$M_N(\alpha)$ is the worst conditioning of the Gaussian linear map $A$ restricted to an
$\alpha$-fraction of coordinates in domain and range. If $A$ is the weight matrix of a one-layer
network with a hard-threshold gate, then $S$ and $T$ are an input and an output activation pattern,
and $M_N(\alpha)$ is the most singular input--output map realized by some pattern; Theorem
\ref{thm:main} gives the rate at which this worst case decays with width. This is the
coordinate-selection counterpart, at speed $N$ and governed by entropy, of the fixed-support hard
edge, which is only polynomial.

\paragraph{Notation.}
$H$ and $h$ are as above (natural logarithm). For a standard $\F$-Gaussian scalar $\zeta$ (so
$\EE|\zeta|^2=1$) and $0<t\le1$ put $p(t)=\PP(|\zeta|\le t)$; then
\begin{equation}\label{eq:pball}
  c_\F\,t^{\dF}\ \le\ p(t)\ \le\ C_\F\,t^{\dF},
\end{equation}
with $p(t)=2\varphi(0)t(1+O(t^2))$ when $\F=\R$ and $p(t)=1-e^{-t^2}$ when $\F=\C$ (there $|\zeta|^2$ is
$\mathrm{Exp}(1)$). More generally, for a standard $\F$-Gaussian vector $g^{(k)}\in\F^k$ put
$p_k(t)=\PP(\|g^{(k)}\|\le t)$; since $\|g^{(k)}\|^2$ is a chi-square variable with $\dF k$ real
degrees of freedom,
\begin{equation}\label{eq:pkball}
  c_{\F,k}\,t^{\dF k}\ \le\ p_k(t)\ \le\ C_{\F,k}\,t^{\dF k},\qquad 0<t\le1,
\end{equation}
and $p_1=p$. For a unit vector $w\in\F^d$ and a standard $\F$-Gaussian vector $g\in\F^d$, the
scalar $\langle g,w\rangle$ is a standard $\F$-Gaussian, by rotation invariance. Constants
$c_\F,C_\F,C_{\F,k}$ depend only on the subscripts and may change between lines.

\section{The first-moment upper bound}\label{sec:upper}
\begin{proposition}\label{prop:upper}
For every $\varepsilon>0$, $\PP\bigl(M_N(\alpha)\le e^{-(2H(\alpha)/\dF+\varepsilon)N}\bigr)\to0$. In
particular $\limsup_N-\tfrac1N\log M_N(\alpha)\le 2H(\alpha)/\dF$ in probability.
\end{proposition}
\begin{proof}
For a single $n\times n$ standard $\F$-Gaussian matrix $G$ one has the hard-edge tail
\begin{equation}\label{eq:hardedge}
  \PP\bigl(\smin(G)\le t\bigr)\ \le\ C_\F\,\bigl(\sqrt n\,t\bigr)^{\dF},\qquad t>0,
\end{equation}
valid for all $t$ (real case: Edelman \cite{edelman} and Rudelson--Vershynin \cite{rv}; complex case
analogously, the smallest squared singular value having a bounded density near $0$ at scale $1/n$).
Each $A[T,S]$ has the law of $G$, so by a union bound over the $\binom Nn^2$ pairs and $\binom
Nn=e^{H(\alpha)N+o(N)}$,
\[
  \PP\bigl(M_N(\alpha)\le t\bigr)\ \le\ \binom Nn^2 C_\F(\sqrt n\,t)^{\dF}
  = e^{2H(\alpha)N+o(N)}\,t^{\dF}.
\]
With $t=e^{-(2H(\alpha)/\dF+\varepsilon)N}$ the right side is $e^{-\dF\varepsilon N+o(N)}\to0$.
\end{proof}

\section{The sharp two-direction small-ball estimate}\label{sec:twoslab}
This is the one new analytic ingredient. Let $u,v$ be independent Haar-random unit vectors in $\F^D$
($D\ge2$) and $g$ a standard $\F$-Gaussian vector in $\F^D$ independent of $u,v$. With $p$ as in
\eqref{eq:pball}, set
\[
  q_D(t)=\PP\bigl(|\langle g,u\rangle|\le t,\ |\langle g,v\rangle|\le t\bigr).
\]
\begin{lemma}\label{lem:twoslab}
There is $C_\F<\infty$, independent of $D$, such that for all $D\ge2$ and $0<t\le1$,
\[
  q_D(t)\ \le\ C_\F\,p(t)^2\,\bigl(1+\log(1/t)\bigr).
\]
\end{lemma}
\begin{proof}
Condition on $u,v$ and write $\rho=\langle u,v\rangle$, $s^2=1-|\rho|^2$. Put $X=\langle
g,u\rangle$, $Y=\langle g,v\rangle$; these are jointly (circularly symmetric) $\F$-Gaussian with unit
variances and $\EE[Y\overline X]=\rho'$, $|\rho'|=|\rho|$. Regressing, $Y=cX+sW$ with $|c|=|\rho|$ and
$W$ a standard $\F$-Gaussian independent of $X$. If $|X|\le t$ then $|cX|\le t$, so
$|sW|=|Y-cX|\le|Y|+|cX|$, whence $\{|X|\le t\}\cap\{|Y|\le t\}\subseteq\{|X|\le t\}\cap\{|W|\le
2t/s\}$. As $X\perp W$,
\begin{equation}\label{eq:cond}
  \PP\bigl(|X|\le t,|Y|\le t\mid u,v\bigr)\ \le\ p(t)\,p(2t/s)\ \le\ C_\F\,p(t)\,\min\!\bigl(1,(t/s)^{\dF}\bigr),
\end{equation}
using $p(r)\le C_\F\min(1,r^{\dF})$ from \eqref{eq:pball}. Now $R:=|\rho|^2=|\langle
u,v\rangle|^2\sim\Beta(a_0,b_0)$ with $(a_0,b_0)=(\tfrac{\dF}2,\tfrac{(D-1)\dF}2)$ (the squared first
coordinate of a Haar unit vector), so $s^2=1-R\sim\Beta(b_0,a_0)$. Taking expectations in
\eqref{eq:cond},
\[
  q_D(t)\ \le\ C_\F\,p(t)\,\EE\bigl[\min(1,(t/s)^{\dF})\bigr],\qquad
  \EE\bigl[\min(1,(t/s)^{\dF})\bigr]=t^{\dF}\,\EE\bigl[s^{-\dF}\one_{\{s\ge t\}}\bigr]+\PP(s<t).
\]
For the first term, with $\alpha_0:=\dF/2$,
\[
  \EE\bigl[s^{-\dF}\one_{\{s\ge t\}}\bigr]=\frac1{B(b_0,a_0)}\int_{t^2}^1 w^{\,b_0-\alpha_0-1}(1-w)^{a_0-1}\,dw,
  \qquad b_0-\alpha_0=\tfrac{(D-2)\dF}2 .
\]
When $D=2$ the exponent $b_0-\alpha_0-1=-1$ and the integral is $\le C_\F(1+\log(1/t))$ (the factor
$(1-w)^{a_0-1}$ is integrable at $w=1$). When $D\ge3$ the exponent exceeds $-1$, the integral
converges, and $B(b_0-\alpha_0,a_0)/B(b_0,a_0)\to1$ and is bounded uniformly in $D$ (as
$\Gamma(x+a)/\Gamma(x)\asymp x^{a}$). Hence $\EE[s^{-\dF}\one_{\{s\ge t\}}]\le C_\F(1+\log(1/t))$
uniformly in $D\ge2$. For the second term, $\PP(s<t)=\PP(\Beta(b_0,a_0)<t^2)\le C_\F\,t^{2b_0}\le
C_\F\,t^{\dF}$ since $2b_0=(D-1)\dF\ge\dF$. Therefore $\EE[\min(1,(t/s)^{\dF})]\le
C_\F\,t^{\dF}(1+\log(1/t))$, and since $t^{\dF}\le p(t)/c_\F$ by \eqref{eq:pball},
\[
  q_D(t)\ \le\ C_\F\,p(t)\cdot t^{\dF}(1+\log(1/t))\ \le\ C_\F'\,p(t)^2\,(1+\log(1/t)).\qedhere
\]
\end{proof}

\begin{remark}\label{rem:crude}
The crude bound $q_D(t)\le p(t)$ (either marginal) loses the factor $p(t)^{-1}=e^{\Theta(N)}$ at the
scales below. Lemma~\ref{lem:twoslab} shows the true excess over independence is only the polynomial
$1+\log(1/t)$, for \emph{every} $D\ge2$; this is what makes the second moment converge at the sharp
constant. The single conditional inequality \eqref{eq:cond} is also the mechanism used for
same-column pairs in Section~\ref{sec:second}, where the overlap is not Haar in a fixed dimension.
\end{remark}

In the second-moment arguments the conditional inequality \eqref{eq:cond} (and its block version
\eqref{eq:vectorslab} below) is averaged over the law of the overlap. The following elementary
estimate packages that step; the factor $\Lambda$ accommodates the polynomial prefactors that the
overlap tails of Section~\ref{sec:second} carry.

\begin{lemma}\label{lem:averaged}
Let $j\ge1$, $q=\dF j$, and let $\sigma^2\in[0,1]$ be a random variable whose lower tail satisfies,
for some $\beta\ge\dF(j+1)/2$, $\gamma\ge1$ and $\Lambda\ge1$,
\[
  \PP\bigl(\sigma^2\le x\bigr)\ \le\ \Lambda\,(\gamma x)^{\beta}\qquad\text{for all }x\in(0,1].
\]
Then, for all $t$ with $\sqrt\gamma\,t\le1$,
\[
  \EE\Bigl[\min\!\bigl(1,(2t/\sigma)^{q}\bigr)\Bigr]\ \le\ C_{\F,j}\,\Lambda\,\gamma^{q/2}\,t^{q}.
\]
\end{lemma}

\begin{proof}
Since $\gamma x\le1$ on the range where the bound is below one, the hypothesis also holds with
$\beta$ replaced by $\beta'=\dF(j+1)/2\le\beta$, so assume $\beta=\beta'$. Write
$\EE[\min(1,(2t/\sigma)^q)]\le\PP(\sigma^2\le4t^2)+(2t)^q\,\EE[\sigma^{-q}\one_{\{\sigma^2>4t^2\}}]$.
The first term is at most $\Lambda(4\gamma t^2)^{\beta'}\le C\Lambda\gamma^{q/2}t^{q}\,
(\gamma^{1/2}t)^{\dF}\le C\Lambda\gamma^{q/2}t^q$. For the second, integrating by parts,
\[
  \EE[\sigma^{-q}\one_{\{\sigma^2>4t^2\}}]\ \le\ 1+\tfrac q2\int_{0}^{1}w^{-q/2-1}\,
  \min\!\bigl(1,\Lambda(\gamma w)^{\beta'}\bigr)\,dw .
\]
Split at $w_*=(\gamma\Lambda^{1/\beta'})^{-1}$: on $(0,w_*]$ the integrand is
$\Lambda\gamma^{\beta'}w^{\beta'-q/2-1}$ and $\beta'-q/2=\dF/2$, giving
$C\Lambda\gamma^{\beta'}w_*^{\dF/2}\le C\Lambda\gamma^{q/2}$; on $[w_*,1]$ the integrand is at most
$w^{-q/2-1}$, giving $Cw_*^{-q/2}\le C\Lambda\gamma^{q/2}$. Multiplying by $(2t)^q$ finishes.
\end{proof}

\section{The marked-row construction}\label{sec:construct}
Fix $\varepsilon\in(0,H(\alpha)/\dF)$ and set $t=t_N=e^{-(2H(\alpha)/\dF-\varepsilon)N}$. Fix a
\emph{marked set} $M\subseteq[N]$ with $|M|=m_N:=\lceil\sqrt N\rceil$; the only properties used are
$m_N\to\infty$ and $m_N=o(N)$. A \emph{configuration} is a triple $c=(S,R,y)$ with $S\subseteq[N]$,
$|S|=n$; $y\in M$; and $R\subseteq[N]\setminus M$, $|R|=n-1$ (so $R$ avoids every marked index, in
particular $y\notin R$). Since $A$ has independent Gaussian entries, the vectors
$\{A_i|_S:i\in R\}$ (rows of $A$ restricted to $S$) are almost surely linearly independent and span a
hyperplane in $\F^S\cong\F^n$; let $u_c\in\F^S$ be a unit normal (defined up to a unit scalar in the
complex case; all events below are invariant under this). Put
\[
  E_c=\bigl\{\,|\langle A_y|_S,\,u_c\rangle|\le t\,\bigr\},\qquad W=W_N=\sum_c \one_{E_c}.
\]
\begin{lemma}\label{lem:suff}
If $W>0$ then $M_N(\alpha)\le t$.
\end{lemma}
\begin{proof}
If $E_c$ holds, set $T=R\cup\{y\}$, so $|T|=n$. For $i\in R$, $\langle A_i|_S,u_c\rangle=0$; hence
$\|A[T,S]\,u_c\|^2=\sum_{i\in T}|\langle A_i|_S,u_c\rangle|^2=|\langle A_y|_S,u_c\rangle|^2\le t^2$.
As $u_c$ is a unit vector, $\smin(A[T,S])\le\|A[T,S]u_c\|\le t$.
\end{proof}
\begin{lemma}\label{lem:mean}
$\EE W=m_N\binom Nn\binom{N-m_N}{n-1}p(t)=\exp(2H(\alpha)N+o(N))\,t^{\dF}=e^{\dF\varepsilon
N+o(N)}\to\infty.$
\end{lemma}
\begin{proof}
Because $y\notin R$, the row $A_y$ is independent of $u_c$, so conditionally on $u_c$ the scalar
$\langle A_y|_S,u_c\rangle$ is a standard $\F$-Gaussian and $\PP(E_c)=p(t)$. The number of
configurations is $m_N\binom Nn\binom{N-m_N}{n-1}$. Since $m_N=o(N)$, $\tfrac1N\log\binom
Nn\to H(\alpha)$ and $\tfrac1N\log\binom{N-m_N}{n-1}\to H(\alpha)$ as well, while $\tfrac1N\log
m_N\to0$; with \eqref{eq:pball} this gives the stated asymptotics.
\end{proof}

\section{The second moment and the lower bound}\label{sec:second}
\begin{proposition}\label{prop:lower}
For every $\varepsilon>0$, $\PP\bigl(M_N(\alpha)>e^{-(2H(\alpha)/\dF-\varepsilon)N}\bigr)\to0$; hence
$\liminf_N-\tfrac1N\log M_N(\alpha)\ge 2H(\alpha)/\dF$ in probability.
\end{proposition}
Theorem~\ref{thm:main} follows from Propositions~\ref{prop:upper} and~\ref{prop:lower}. By
Lemma~\ref{lem:suff} and Chebyshev, $\PP(M_N>t)\le\PP(W=0)\le\Var(W)/(\EE W)^2$, so it suffices to
show $\Var(W)=o((\EE W)^2)$. Write $B=\binom Nn\binom{N-m_N}{n-1}=e^{2H(\alpha)N+o(N)}$, so $\EE
W=m_N\,B\,p(t)$ and $(\EE W)^2=m_N^2B^2p(t)^2$. We must control $\EE[W^2]=\sum_{c,c'}\PP(E_c\cap
E_{c'})$.

\subsection{Distinct marked rows give independence}
Let $c=(S,R,y)$, $c'=(S',R',y')$ and $\mathcal G=\sigma(A_i:i\neq y)$. The marked set is the point
of the restriction: since $y,y'\in M$ while $R,R'\subseteq[N]\setminus M$, we always have $y'\notin
R$ and $y\notin R'$, and also $y\neq y'\Rightarrow y\notin\{y'\}$.
\begin{lemma}\label{lem:indep}
If $y\neq y'$ then $\PP(E_c\cap E_{c'})=p(t)^2$.
\end{lemma}
\begin{proof}
The row $A_y$ enters neither $u_c$ (built from $R\not\ni y$), nor $u_{c'}$ (from $R'\not\ni y$), nor
the marked row $A_{y'}$ of $c'$ (as $y'\neq y$); hence both normals and the event $E_{c'}$ are
$\mathcal G$-measurable, while $E_c=\{|\langle A_y|_S,u_c\rangle|\le t\}$ depends on the fresh row
$A_y$ with $\PP(E_c\mid\mathcal G)=p(t)$ for every value of $u_c$. Since $A_y\perp\mathcal G$,
$\PP(E_c\cap E_{c'})=\EE[\one_{E_{c'}}\PP(E_c\mid\mathcal G)]=p(t)\,\PP(E_{c'})=p(t)^2$.
\end{proof}
Thus the only correlated pairs have $y=y'$: a marked row can never be another configuration's
normal-defining row, which is exactly what the restriction $R\subseteq[N]\setminus M$ buys.

\subsection{Same marked row}\label{sub:same}
Here $y=y'$ and $A_y=A_{y'}$ is a single fresh row. Given $\mathcal G$, $u_c,u_{c'}$ are determined,
and $(\langle A_y|_S,u_c\rangle,\langle A_y|_{S'},u_{c'}\rangle)=(\langle A_y,\tilde
u_c\rangle,\langle A_y,\tilde u_{c'}\rangle)$, where $\tilde u_c,\tilde u_{c'}\in\F^N$ are the
zero-padded normals. By the conditional inequality \eqref{eq:cond} (applied to the standard Gaussian
$A_y$ and the unit vectors $\tilde u_c,\tilde u_{c'}$, with $R:=|\langle\tilde u_c,\tilde
u_{c'}\rangle|^2$ and $s^2=1-R$),
\begin{equation}\label{eq:sameG}
  \PP(E_c\cap E_{c'}\mid\mathcal G)\ \le\ C_\F\,p(t)\,\min\!\bigl(1,(t/s)^{\dF}\bigr).
\end{equation}
The overlap $R=|\langle\tilde u_c,\tilde u_{c'}\rangle|^2$ is controlled by two tail bounds, one
through the columns and one through the rows. Write $D_{\mathrm c}=n-|S\cap S'|$ and $D_{\mathrm
r}=n-|R\cap R'|$, so $D_{\mathrm r}\ge1$ with equality iff $R=R'$.

\begin{lemma}\label{lem:overlaptail}
For two configurations $c,c'$ (with any marked data) and all $x\in(0,1]$:
\begin{enumerate}
\item (row bound) if $D_{\mathrm r}\ge2$, then
  $\PP\bigl(R>1-x\bigr)\le C_\F\,x^{(D_{\mathrm r}-1)\dF/2}$, with $C_\F$ uniform in $D_{\mathrm r}$
  and $N$;
\item (column bound) if $D_{\mathrm c}\ge1$, then
  $\PP\bigl(R>1-x\bigr)\le n\,\bigl(C_\F\,n\,x/D_{\mathrm c}\bigr)^{D_{\mathrm c}\dF/2}$.
\end{enumerate}
\end{lemma}

\begin{proof}
(ii) Since $\langle\tilde u_c,\tilde u_{c'}\rangle=\langle u_c|_{S\cap S'},u_{c'}|_{S\cap
S'}\rangle$ and $u_{c'}$ is a unit vector, $R\le\|u_c|_{S\cap S'}\|^2=1-Z$ with
$Z=\|u_c|_{S\setminus S'}\|^2$. The normal $u_c$ is Haar on the unit sphere of $\F^S$ (the span of
$n-1$ i.i.d.\ Gaussian vectors is a uniformly random hyperplane), so
$Z\sim\Beta(a,b)$ with $a=D_{\mathrm c}\dF/2$, $b=(n-D_{\mathrm c})\dF/2$. We may assume $x\le1/2$
and $b\ge1$, since otherwise the stated bound exceeds one. Hence, bounding
$(1-w)^{b-1}\le1$,
\[
  \PP(Z\le x)\ \le\ \frac{x^a}a\cdot\frac1{B(a,b)}
  =x^a\,\frac{\Gamma(a+b)}{\Gamma(a+1)\Gamma(b)}
  \ \le\ x^a\,(a+b)^{a+1}\Bigl(\frac ea\Bigr)^a
  \ \le\ n\,\Bigl(\frac{C_\F nx}{D_{\mathrm c}}\Bigr)^{a},
\]
using $\Gamma(a+b)\le\Gamma(b)(a+b)^{a+1}$ and $\Gamma(a+1)\ge(a/e)^a$.

(i) Condition on all rows indexed by $R$; this determines $u_c$. Set
$V=\mathrm{span}\{A_i|_{S'}:i\in R\cap R'\}\subseteq\F^{S'}$, of dimension $n-D_{\mathrm r}$ almost
surely, and $W_0=V^\perp\subseteq\F^{S'}$, of dimension $D_{\mathrm r}$. The $D_{\mathrm r}-1$ rows
indexed by $R'\setminus R$ are independent of the conditioning, their restrictions to $S'$ are
standard Gaussians in $\F^{S'}$, and their orthogonal projections onto $W_0$ are independent
standard Gaussians in $W_0$. The normal $u_{c'}$ lies in $W_0$ and is orthogonal to these
projections, so, given the conditioning, $u_{c'}$ is Haar on the unit sphere of $W_0$. If
$\hat w$ denotes the unit vector along the projection of $\tilde u_c$ onto $W_0$ (when the
projection vanishes, $R=0$), then
$R=|\langle\tilde u_c,u_{c'}\rangle|^2\le|\langle\hat w,u_{c'}\rangle|^2=:1-Z'$, and
$Z'\sim\Beta(b_0,a_0)$ with $b_0=(D_{\mathrm r}-1)\dF/2$, $a_0=\dF/2$. Hence
\[
  \PP(Z'\le x)\ \le\ \frac{x^{b_0}}{b_0}\cdot\frac1{B(b_0,a_0)}
  =x^{b_0}\,\frac{\Gamma(a_0+b_0)}{\Gamma(b_0+1)\Gamma(a_0)}\ \le\ C_\F\,x^{b_0},
\]
because $a_0\le1$ gives $\Gamma(a_0+b_0)\le\max\{\Gamma(b_0),\Gamma(b_0+1)\}\le
C\,\Gamma(b_0+1)$ by the log-convexity of $\Gamma$. Averaging over the conditioning preserves both
bounds.
\end{proof}

We now bound $\EE[\one_{E_c\cap E_{c'}}]=\EE[\PP(E_c\cap E_{c'}\mid\mathcal G)]$ over three regions
of $(D_{\mathrm c},D_{\mathrm r})$. Recall from \eqref{eq:sameG} that the quantity to average is
$C_\F\,p(t)\min(1,(t/s)^{\dF})$ with $s^2=1-R$; Lemma~\ref{lem:averaged} (with $j=1$, $q=\dF$)
converts each tail bound of Lemma~\ref{lem:overlaptail} into a bound on this average. Note
$\sqrt\gamma\,t\le1$ holds in all cases below since $\gamma\le C_\F n$ and $t=e^{-\Theta(N)}$. Two
counting facts are used: the fraction of ordered pairs $(S,S')$ with $D_{\mathrm c}\le a$ is at most
$N^{2a}\binom Nn^{-1}=N^{2a}e^{-H(\alpha)N+o(N)}$, and likewise for rows.

\emph{Region $D_{\mathrm r}\ge3$.} The row bound gives the hypothesis of Lemma~\ref{lem:averaged}
with $\beta=(D_{\mathrm r}-1)\dF/2\ge\dF$, $\gamma=C_\F$, $\Lambda=1$, so
$\PP(E_c\cap E_{c'})\le C_\F\,p(t)\cdot t^{\dF}\le C_\F\,p(t)^2$. Summed over the $m_N$ choices of
$y=y'$ and all such pairs, the contribution is at most $C_\F(\EE W)^2/m_N=o((\EE W)^2)$.

\emph{Region $D_{\mathrm r}\le2$, $D_{\mathrm c}\ge2$.} The column bound gives the hypothesis with
$\beta=D_{\mathrm c}\dF/2\ge\dF$, $\gamma=C_\F n/D_{\mathrm c}$, $\Lambda=n$, so $\PP(E_c\cap
E_{c'})\le C_\F\,n\,(n/D_{\mathrm c})^{\dF/2}\,p(t)^2\le N^{C}p(t)^2$. The rows are pinned:
the fraction of pairs with $D_{\mathrm r}\le2$ is at most $N^{4}e^{-H(\alpha)N+o(N)}$, so the
contribution is at most $N^{C}e^{-H(\alpha)N+o(N)}(\EE W)^2=o((\EE W)^2)$.

\emph{Region $D_{\mathrm r}\le2$, $D_{\mathrm c}\le1$.} Use the crude bound $\PP(E_c\cap
E_{c'}\mid\mathcal G)\le p(t)$ (one marginal), so the ratio to $p(t)^2$ is at most
$C_\F\,t^{-\dF}=C_\F e^{(2H(\alpha)-\dF\varepsilon)N}$. Both coordinates are pinned: the fraction of
such pairs is at most $N^{C}e^{-2H(\alpha)N+o(N)}$. The contribution is at most
$N^{C}e^{-\dF\varepsilon N+o(N)}(\EE W)^2=o((\EE W)^2)$; this region includes the diagonal $c=c'$,
whose contribution $\EE W$ is $o((\EE W)^2)$ as well.

\subsection{Conclusion}
By Lemma~\ref{lem:indep} the distinct-marked-row pairs contribute $\le(\#\text{configs})^2p(t)^2=(\EE
W)^2$, and the same-marked-row pairs, split into the three regions above, contribute $o((\EE W)^2)$.
Hence $\EE[W^2]\le(\EE W)^2+o((\EE W)^2)$, i.e.\ $\Var(W)=o((\EE W)^2)$. By Chebyshev $\PP(W=0)\to0$,
and Lemma~\ref{lem:suff} yields Proposition~\ref{prop:lower}. \qed

\begin{remark}
The marked set $M$ is what makes the second moment converge \emph{at the sharp constant}. Without it,
pairs in which each configuration's marked row is the other's normal-defining row form a $\Theta(1)$
fraction of all pairs and are genuinely correlated, so no crude bound controls them; forbidding
normal rows from $M$ removes them, leaving only the same-marked-row pairs, whose $\Theta(1/m_N)$
weight and the estimates of \S\ref{sub:same} suffice. Any $m_N$ with $m_N\to\infty$, $m_N=o(N)$
works; $m_N=\lceil\sqrt N\rceil$ is one choice.
\end{remark}

\section{Surplus rows}\label{sec:surplus}
Fix an integer $r\ge1$ and write $k=r+1$; throughout this section $r$ (hence $k$) is fixed as
$N\to\infty$, and constants may depend on $\F$ and $k$. We prove Theorem~\ref{thm:surplus}. The
architecture is that of Sections~\ref{sec:construct}--\ref{sec:second} with $k$ marked rows per
configuration instead of one; the new points are a vector version of the two-direction estimate, a
second tail bound for the overlap of two constructed normals, and a bookkeeping of pairs by the
number of shared marked rows.

\subsection{Upper bound}
For a single $(n+r)\times n$ standard $\F$-Gaussian matrix $G$ the rectangular hard-edge tail
\begin{equation}\label{eq:recttail}
  \PP\bigl(\smin(G)\le t\bigr)\ \le\ \bigl(C_{\F,r}\,\sqrt n\,t\bigr)^{\dF(r+1)},\qquad t>0,
\end{equation}
holds \cite[Lemma 2.2]{shmalo} (see also \cite{rv}). A union bound over the
$\binom Nn\binom N{n+r}=e^{2H(\alpha)N+o(N)}$ support pairs gives, for
$t=e^{-(2H(\alpha)/(\dF k)+\varepsilon)N}$,
\[
  \PP\bigl(M_N^{(r)}(\alpha)\le t\bigr)\ \le\ e^{2H(\alpha)N+o(N)}\,t^{\dF k}
  \ =\ e^{-\dF k\varepsilon N+o(N)}\ \to\ 0 .
\]

\subsection{Construction and mean}
Fix $\varepsilon\in(0,H(\alpha)/(\dF k))$ and set $t=t_N=e^{-(2H(\alpha)/(\dF k)-\varepsilon)N}$.
Keep the marked set $M\subseteq[N]$, $|M|=m_N=\lceil\sqrt N\rceil$. A configuration is now a triple
$c=(S,R,Y)$ with $|S|=n$; $Y\subseteq M$, $|Y|=k$; and $R\subseteq[N]\setminus M$, $|R|=n-1$. As
before $u_c\in\F^S$ is the unit normal to $\mathrm{span}\{A_i|_S:i\in R\}$, and
\[
  E_c=\Bigl\{\,\sum_{y\in Y}|\langle A_y|_S,\,u_c\rangle|^2\le t^2\,\Bigr\},\qquad
  W=\sum_c\one_{E_c}.
\]
If $E_c$ holds, then $T=R\cup Y$ has $|T|=n+r$ and $\|A[T,S]u_c\|^2=\sum_{y\in Y}|\langle
A_y|_S,u_c\rangle|^2\le t^2$, so $\smin(A[T,S])\le t$ and $M_N^{(r)}\le t$. Since the rows indexed
by $Y$ are independent of $u_c$, the vector $(\langle A_y|_S,u_c\rangle)_{y\in Y}$ is standard
$\F$-Gaussian in $\F^k$ given $u_c$, so $\PP(E_c)=p_k(t)$, and with
$\binom{m_N}k=e^{o(N)}$,
\[
  \EE W=\binom Nn\binom{m_N}k\binom{N-m_N}{n-1}\,p_k(t)
  =e^{2H(\alpha)N+o(N)}\,t^{\dF k}=e^{\dF k\varepsilon N+o(N)}\ \to\ \infty .
\]

\subsection{A vector two-direction estimate}
The two-direction estimate of Section~\ref{sec:twoslab} extends to blocks of $j$ rows with no change
in the proof. Let $u,v$ be unit vectors in $\F^d$ with $|\langle u,v\rangle|^2=\omega$ and
$s^2=1-\omega$, and let $g_1,\dots,g_j$ be independent standard $\F$-Gaussian vectors in $\F^d$.
Set $X=(\langle g_i,u\rangle)_{i\le j}$ and $Y=(\langle g_i,v\rangle)_{i\le j}$, standard
$\F$-Gaussian vectors in $\F^j$ with componentwise regression $Y=cX+sW$, $|c|^2=\omega$, where
$W\in\F^j$ is standard and independent of $X$ (each pair $(X_i,Y_i)$ is jointly Gaussian with the
same correlation $c$, and the pairs are independent). On $\{\|X\|\le t\}$ one has $\|sW\|\le\|Y\|+
|c|\,\|X\|\le2t$ whenever $\|Y\|\le t$, so
\begin{equation}\label{eq:vectorslab}
  \PP\bigl(\|X\|\le t,\ \|Y\|\le t\ \big|\ u,v\bigr)\ \le\ p_j(t)\,p_j(2t/s)
  \ \le\ C_{\F,j}\,p_j(t)\,\min\!\bigl(1,(t/s)^{\dF j}\bigr).
\end{equation}

The averaged form is supplied by Lemma~\ref{lem:averaged} with $q=\dF j$, fed by the overlap tails
of Lemma~\ref{lem:overlaptail}. That lemma was stated and proved for two arbitrary configurations;
its hypotheses involve only the column sets and the normal-row sets, so it applies verbatim to the
configurations of this section, whose normals are built from $R,R'$ exactly as before. As there, we
write $\omega=|\langle\tilde u_c,\tilde u_{c'}\rangle|^2$ for the squared overlap of the zero-padded
normals, $s^2=1-\omega$, and $D_{\mathrm c}=n-|S\cap S'|$, $D_{\mathrm r}=n-|R\cap R'|$.

\subsection{Second moment}
Let $c,c'$ be two configurations and $j=|Y\cap Y'|\in\{0,\dots,k\}$. As in
Section~\ref{sec:second} let $\mathcal G=\sigma(A_i:i\notin Y\cup Y')$; both normals are $\mathcal
G$-measurable because $R,R'$ avoid $M\supseteq Y\cup Y'$.

\emph{Disjoint marked sets ($j=0$).} Given $\mathcal G$, the event $E_c$ depends only on the rows
in $Y$ and $E_{c'}$ only on the rows in $Y'$, disjoint families of fresh rows, and each conditional
probability equals $p_k(t)$ exactly. Hence $\PP(E_c\cap E_{c'})=p_k(t)^2$, and these pairs
contribute at most $(\EE W)^2$.

\emph{Shared marked rows ($1\le j\le k$).} The fraction of ordered configuration pairs with
$|Y\cap Y'|=j$ is $\binom kj\binom{m_N-k}{k-j}/\binom{m_N}k\le C_k\,m_N^{-j}$. Split each event
along shared and own marked rows:
\[
  E_c\ \subseteq\ \Bigl\{\sum_{y\in Y\cap Y'}|\langle A_y,\tilde u_c\rangle|^2\le t^2\Bigr\}
  \cap\Bigl\{\sum_{y\in Y\setminus Y'}|\langle A_y,\tilde u_c\rangle|^2\le t^2\Bigr\},
\]
and similarly for $E_{c'}$. Given $\mathcal G$, the three row families $Y\cap Y'$, $Y\setminus Y'$,
$Y'\setminus Y$ are independent, the own-row events have probabilities $p_{k-j}(t)$ each, and the
joint shared-row event is bounded by \eqref{eq:vectorslab} with $d=N$, $u=\tilde u_c$,
$v=\tilde u_{c'}$:
\[
  \PP(E_c\cap E_{c'}\mid\mathcal G)\ \le\ p_{k-j}(t)^2\cdot C\,p_j(t)\,
  \min\!\bigl(1,(t/s)^{\dF j}\bigr),\qquad s^2=1-\omega .
\]
Taking expectations, using Lemma~\ref{lem:overlaptail} through Lemma~\ref{lem:averaged}
($q=\dF j$), and noting $p_{k-j}(t)^2p_j(t)\,t^{\dF j}\le C\,p_k(t)^2$ by \eqref{eq:pkball}, the
pairs with a given $j\ge1$ split into three regions, exactly as in Section~\ref{sec:second}. Recall
$t^{-\dF k}=e^{(2H(\alpha)-\dF k\varepsilon)N}$, and that the fraction of ordered column pairs with
$D_{\mathrm c}\le a$ is at most $N^{2a}e^{-H(\alpha)N+o(N)}$, likewise for rows.

\emph{Region $D_{\mathrm r}\ge j+2$.} The row bound of Lemma~\ref{lem:overlaptail} verifies the
hypothesis of Lemma~\ref{lem:averaged} with $\beta=(D_{\mathrm r}-1)\dF/2\ge\dF(j+1)/2$,
$\gamma=C_\F$, $\Lambda=1$, so $\PP(E_c\cap E_{c'})\le C\,p_{k-j}(t)^2p_j(t)\,t^{\dF j}\le
C\,p_k(t)^2$, uniformly over the region. With the factor $m_N^{-j}$ from the shared marked rows the
contribution is at most $C\,m_N^{-j}(\EE W)^2=o((\EE W)^2)$.

\emph{Region $D_{\mathrm r}\le j+1$, $D_{\mathrm c}\ge j+1$.} The column bound verifies the
hypothesis with $\beta=D_{\mathrm c}\dF/2\ge\dF(j+1)/2$, $\gamma=C_\F n/D_{\mathrm c}$,
$\Lambda=n$, so $\PP(E_c\cap E_{c'})\le C\,n\,(n/D_{\mathrm c})^{\dF j/2}p_k(t)^2\le
N^{C_k}p_k(t)^2$. The rows are pinned ($D_{\mathrm r}\le j+1\le k+1$), so the pair fraction is at
most $N^{C_k}e^{-H(\alpha)N+o(N)}$, and the contribution is $o((\EE W)^2)$.

\emph{Region $D_{\mathrm r}\le j+1$, $D_{\mathrm c}\le j$.} Use the crude bound: dropping the
$E_{c'}$ shared-row constraint, $\PP(E_c\cap E_{c'}\mid\mathcal G)\le p_j(t)\,p_{k-j}(t)^2$, so the
ratio to $p_k(t)^2$ is at most $C\,t^{-\dF j}\le C\,t^{-\dF k}=Ce^{(2H(\alpha)-\dF
k\varepsilon)N}$. Both coordinates are pinned, so the pair fraction is at most
$N^{C_k}e^{-2H(\alpha)N+o(N)}$, and the contribution is at most $N^{C_k}e^{-\dF k\varepsilon
N+o(N)}(\EE W)^2=o((\EE W)^2)$; this region includes the diagonal, whose contribution is $\EE W$.

Summing over $j$, $\EE[W^2]\le(\EE W)^2(1+o(1))$, so $\Var(W)=o((\EE W)^2)$, and Chebyshev's
inequality gives $W>0$ with probability tending to one, hence $M_N^{(r)}\le t$. Combined with the
upper bound this proves Theorem~\ref{thm:surplus}. \qed

\begin{remark}
The row bound of Lemma~\ref{lem:overlaptail} is what makes the $k\ge2$ case work: the column bound
alone leaves the pairs with small $D_{\mathrm c}$ but generic $D_{\mathrm r}$ uncontrolled, and
there the row bound takes over --- two normals built on row sets that differ in many places cannot
be nearly parallel, no matter how similar their column sets are. Only pairs pinned in both
coordinates are genuinely correlated, and those are exponentially rare.
\end{remark}

\begin{remark}\label{rem:moderate}
The argument is uniform in a slowly growing surplus. For $r=r_N\to\infty$ the only steps that see
$r$ are the shared-marked-row count, which requires $\sum_{j\ge1}\binom kj^2m_N^{-j}C^j\to0$ and
hence $r^2=o(m_N)$, and various factors $e^{O(r\log N)}$ (binomials in $k$, the constants
$C_{\F,k}$ of \eqref{eq:pkball}, the pinned-pair counts $N^{C_k}$), which are absorbed because the
comparison budget $e^{\dF k\varepsilon N}$ itself grows linearly in $k$. With $m_N=\lceil\sqrt
N\rceil$ this gives Theorem~\ref{thm:surplus} uniformly for $r=o(N^{1/4})$, in the form
$-\tfrac{r+1}N\log M_N^{(r)}(\alpha)\to2H(\alpha)/\dF$. We have not pursued the optimal range of
$r_N$.
\end{remark}

\section{Consequences and complements}\label{sec:conseq}

\subsection{Support-optimized products}
Let $\kappa\ge1$ and let $A_1,\dots,A_\kappa$ be independent $N\times N$ standard $\F$-Gaussian
matrices. For supports $S_0,\dots,S_\kappa\subseteq[N]$ of size $n$ let
$M(\mathbf S)=P_{S_\kappa}A_\kappa\cdots P_{S_1}A_1P_{S_0}$, viewed as a map $\F^{S_0}\to
\F^{S_\kappa}$, where $P_S$ is the coordinate restriction.

\begin{corollary}\label{cor:products}
In probability,
\[
  \liminf_{N\to\infty}\ -\frac1N\log\ \min_{\mathbf S}\ \smin\bigl(M(\mathbf S)\bigr)
  \ \ge\ \frac{2H(\alpha)}{\dF}.
\]
\end{corollary}

\begin{proof}
By Theorem~\ref{thm:main} applied to $A_1$, with probability tending to one there are $S_0,S_1$ and
a unit $v\in\F^{S_0}$ with $\|P_{S_1}A_1P_{S_0}v\|\le e^{-(2H(\alpha)/\dF-\varepsilon)N}$. Choose
the remaining supports arbitrarily. Then
\[
  \smin\bigl(M(\mathbf S)\bigr)\le\|M(\mathbf S)v\|
  \le\prod_{i=2}^{\kappa}\|A_i\|\cdot\|P_{S_1}A_1P_{S_0}v\|
  \le(3\sqrt N)^{\kappa-1}e^{-(2H(\alpha)/\dF-\varepsilon)N}
\]
with probability tending to one, using the standard bound $\|A_i\|\le3\sqrt N$. The polynomial
factor is $e^{o(N)}$.
\end{proof}

For $\kappa=1$ this is sharp by Theorem~\ref{thm:main}. For $\kappa\ge2$ the best available upper
bound remains the union bound over all $\kappa+1$ supports, with exponent
$(\kappa+1)H(\alpha)/\dF$, and the sharp constant is open; the natural guess is that one factor
carries all the singularity and the exponent stays $2H(\alpha)/\dF$. One piece of evidence can be
made precise: selecting the \emph{output} coordinates of a factor cannot shrink a fixed direction
at speed $N$.

\begin{lemma}\label{lem:outputselect}
Let $A$ be $N\times N$ standard $\F$-Gaussian and $u\in\F^N$ a fixed unit vector. For every $b>0$,
\[
  \PP\Bigl(\min_{|S|=n}\bigl\|P_SAu\bigr\|\le e^{-bN}\Bigr)\ \le\
  e^{H(\alpha)N}\bigl(C_\F e^{-bN}\bigr)^{\dF n}\ =\ e^{-\dF\alpha bN^2+O(N)} .
\]
\end{lemma}

\begin{proof}
The coordinates of $Au$ are i.i.d.\ standard $\F$-Gaussians $\zeta_1,\dots,\zeta_N$. If some $S$
with $|S|=n$ has $\|P_SAu\|\le e^{-bN}$, then at least $n$ of the $\zeta_i$ satisfy
$|\zeta_i|\le e^{-bN}$, an event of probability at most
$\binom Nn\,p(e^{-bN})^{n}\le e^{H(\alpha)N}(C_\F e^{-bN})^{\dF n}$.
\end{proof}

Thus in a product only the two supports adjacent to a single factor can buy singularity at speed
$N$ through that factor, which is what the guess $2H(\alpha)/\dF$ expresses; what remains open is
to control the exponentially many composed input directions that a chain of factors makes
available.

\subsection{Rademacher entries}
The constant in Theorem~\ref{thm:main} is a Gaussian (more precisely, absolutely continuous)
phenomenon; for discrete entries the minimum can vanish outright.

\begin{proposition}\label{prop:rademacher}
Let $B$ be $N\times N$ with independent uniform $\pm1$ entries and let $M_N^{\pm}(\alpha)$ be as in
\eqref{eq:MN} with $A$ replaced by $B$.
\begin{enumerate}
\item If $\alpha\le\tfrac12$, then $M_N^{\pm}(\alpha)=0$ with probability tending to one.
\item If $2H(\alpha)<\alpha\log2$ (numerically $\alpha>\alpha_0$, $\alpha_0=0.9051\ldots$ the root
  of $2H(\alpha)=\alpha\log2$), then $M_N^{\pm}(\alpha)>0$ with probability tending to one.
\end{enumerate}
\end{proposition}

\begin{proof}
(i) Pair up the columns as $(1,2),(3,4),\dots$; for the $i$-th pair let $Z_i$ be the number of rows
on which the two columns agree. The $Z_i$ are independent $\mathrm{Bin}(N,\tfrac12)$ variables, so
$\PP(Z_i\ge\lfloor N/2\rfloor)\ge\tfrac12$ for each $i$, and with probability at least
$1-2^{-\lfloor N/2\rfloor}$ some pair agrees on at least $\lfloor N/2\rfloor\ge n$ rows. Taking $T$
inside the agreement set of that pair and $S$ containing the two columns produces a submatrix with
two equal columns.

(ii) By \cite{tikhomirov}, $\PP(\text{an $n\times n$ $\pm1$ matrix is singular})=(\tfrac12+o(1))^n$.
A union bound over the $\binom Nn^2$ submatrices gives
$\PP(\exists\,T,S:\ B[T,S]\text{ singular})\le e^{(2H(\alpha)-\alpha\log2)N+o(N)}\to0$.
\end{proof}

For $\alpha\in(\tfrac12,\alpha_0)$ the expected number of singular submatrices diverges and we
expect $M_N^{\pm}=0$ there too; the second-moment question is open. In particular no
distribution-free version of Theorem~\ref{thm:main} can hold below $\alpha_0$, and the correct
universality class (continuous entry laws, say) is an open question.

\subsection{Tall submatrices}
For $\beta>\alpha$ the behaviour is qualitatively different. A single $(\beta N)\times(\alpha N)$
Gaussian submatrix is well conditioned, $\smin\asymp(\sqrt\beta-\sqrt\alpha)\sqrt N$ by Bai--Yin, and
support optimization cannot spoil this: the minimum stays of order $\sqrt N$.

\begin{proposition}\label{prop:tall}
Fix $0<\alpha<\beta<1$. There is $c_1=c_1(\alpha,\beta,\F)>0$ such that, in probability,
\[
  c_1\sqrt N\ \le\ \min_{|T|=\lfloor\beta N\rfloor,\ |S|=\lfloor\alpha N\rfloor}\smin\bigl(A[T,S]\bigr)
  \ \le\ \bigl(\sqrt\beta-\sqrt\alpha+o(1)\bigr)\sqrt N .
\]
\end{proposition}

\begin{proof}
Put $m=\lfloor\beta N\rfloor$, $n=\lfloor\alpha N\rfloor$. The upper bound is Bai--Yin for one fixed
submatrix. For the lower bound fix a support $(T,S)$ and a level $c\sqrt N$ with $c\le\sqrt\beta/2$.
For a fixed unit $u\in\F^S$, $\|A[T,S]u\|^2=\sum_{i\in T}|\langle A_i|_S,u\rangle|^2$ is a sum of $m$
independent squared standard $\F$-Gaussians (mean $m$), so a chi-square lower tail gives
$\PP(\|A[T,S]u\|\le2c\sqrt N)\le(C_\F c^2/\beta)^{\dF m/2}$. Set $\tau=\sqrt{2(H(\beta)+H(\alpha))+2}$
and let $\mathcal N$ be a $\rho$-net of the unit sphere of $\F^S$ with $|\mathcal N|\le(3/\rho)^{\dF
n}$ and $\rho=c/(\sqrt\beta+\sqrt\alpha+\tau)$. On the event
$\|A[T,S]\|\le(\sqrt\beta+\sqrt\alpha+\tau)\sqrt N$, whose complement has probability at most
$e^{-\tau^2N/2+o(N)}$, any unit $u$ with $\smin(A[T,S])\le c\sqrt N$ realized at $u$ has a net point
$u_0$ with $\|A[T,S]u_0\|\le c\sqrt N+\|A[T,S]\|\rho\le2c\sqrt N$. Hence
\[
  \PP\bigl(\smin(A[T,S])\le c\sqrt N\bigr)\ \le\
  \Bigl(\tfrac{3(\sqrt\beta+\sqrt\alpha+\tau)}{c}\Bigr)^{\dF n}\bigl(C_\F c^2/\beta\bigr)^{\dF m/2}
  +e^{-\tau^2N/2+o(N)}.
\]
Summing over the $\binom Nm\binom Nn=e^{(H(\beta)+H(\alpha))N+o(N)}$ supports, the second term
contributes $e^{-N+o(N)}$ (as $\tau^2/2=H(\beta)+H(\alpha)+1$), and the first contributes
$e^{-\eta N+o(N)}$ with
\[
  \eta=\dF(\beta-\alpha)\log(1/c)-H(\beta)-H(\alpha)-\dF\alpha\log\bigl(3(\sqrt\beta+\sqrt\alpha+\tau)\bigr)
  -\tfrac{\dF\beta}2\log(C_\F/\beta),
\]
which is positive for all small enough $c$ because $\beta>\alpha$. Thus $\PP(\min\le c\sqrt N)\to0$.
\end{proof}

\begin{remark}
The term $\dF(\beta-\alpha)\log(1/c)$ is the hard-edge cost of driving the $(\beta-\alpha)N$ excess
rows into near-dependence; it diverges as $c\to0$, which is why the exponential support entropy
cannot reach zero. As $\beta\downarrow\alpha$ the admissible $c_1$ tends to $0$, matching the square
case, where the entropy overcomes the now $O(1)$ hard-edge exponent and drives the minimum to
$e^{-\Theta(N)}$. The sharp constant $\lim N^{-1/2}\min$ is the solution of
$H(\beta)+H(\alpha)=\dF\,I(c;\alpha,\beta)$, with $I$ the left-tail rate of the smallest singular
value of a rectangular Gaussian ensemble of aspect ratio $\beta/\alpha$; we do not determine it.
\end{remark}

\section{Numerical confirmation}\label{sec:num}
Two checks accompany the proof (scripts in \texttt{numerics/}). First, the two-direction estimate
of Lemma~\ref{lem:twoslab}: Monte Carlo gives $q_D(t)/p(t)^2$ growing like $\log(1/t)$ at $D=2$
(real: $2.2,2.5,2.8$; complex: $3.8,4.5,5.1$ at $t=0.15,0.10,0.06$), decaying to $1$ as $D$ grows,
and everywhere far below the crude ratio $1/p(t)$ — confirming the polynomial (not exponential)
excess over independence. Second, an exhaustive evaluation of $M_N(\alpha)$ at $\alpha=1/2$ and its
comparison with the first-moment scale $t_\ast$ solving $\binom Nn^2 C_\F(\sqrt n\,t_\ast)^{\dF}=1$
(with $C_\F$ calibrated separately): the measured $-\tfrac1N\log M_N$ equals the first-moment
prediction to within the sampling error — real ($N=8,10,12$): $1.19,1.24,1.24$ versus $1.14,1.18,1.21$;
complex ($N=6,8,10$): $0.64,0.64,0.68$ versus $0.59,0.62,0.63$ — and both rise toward $2H(1/2)/\dF$
($=1.386$ real, $0.693$ complex), the slow $O(N^{-1}\log N)$ approach accounting for the finite-$N$
gap. That the minimum sits at the first-moment scale is the finite-$N$ signature of the second-moment
concentration proved above.

The same check was run for one surplus row ($r=1$), where the exponent should approach
$2H(1/2)/(2\dF)$, i.e.\ $0.693$ (real) and $0.347$ (complex), and the first-moment scale $t_\ast$
solves $\binom Nn\binom N{n+1}C_\F(\sqrt n\,t_\ast)^{2\dF}=1$. Measured against predicted: real
($N=8,10,12$): $0.560,0.601,0.598$ versus $0.559,0.590,0.606$; complex ($N=8,10$): $0.311,0.329$
versus $0.305,0.319$. Again the minimum sits at the first-moment scale in both fields
(\texttt{numerics/verify\_surplus\_finite\_size.py}).

\begin{thebibliography}{9}
\bibitem{shmalo} Y.~Shmalo, \emph{Extreme least singular values of Gaussian row submatrices}, 2026.
\bibitem{edelman} A.~Edelman, \emph{Eigenvalues and condition numbers of random matrices}, SIAM J.
  Matrix Anal. Appl. \textbf{9} (1988), 543--560.
\bibitem{rv} M.~Rudelson and R.~Vershynin, \emph{The Littlewood--Offord problem and invertibility of
  random matrices}, Adv. Math. \textbf{218} (2008), 600--633.
\bibitem{tikhomirov} K.~Tikhomirov, \emph{Singularity of random Bernoulli matrices}, Ann. of Math.
  \textbf{191} (2020), 593--634.
\end{thebibliography}
\end{document}
